\label{sec:limitations}

\subsection{Relationship to prior confounder-audit and benchmarking work}

Dawood et al.\ perform a similarly motivated, larger-scale confounder audit but do not include
prostate cancer, do not cross-validate across independently trained encoders, and do not test
whether qualification changes a marker's value for a downstream clinical outcome -- the
comparison we make directly in Section~\ref{sec:leopard-results}. Gevaert/Bareja et al.\ and
Neidlinger et al.\ benchmark many encoders across many tasks, including several of the exact
targets we study (TCGA-PRAD ERG status, SICAPv2 grading), but their reported methodology does
not appear to enforce patient-disjoint evaluation, which our own early results show is
necessary to avoid a specific, previously undetected pseudo-replication failure mode. Our
contribution is narrower than either in absolute scope (prostate-specific, seven markers) but
adds three things neither provides together: mandatory patient-disjoint discipline throughout,
cross-encoder replication of the entire marker pool rather than a subset, and a direct
same-outcome test of whether qualification is consequential rather than merely descriptive.

\subsection{Limitations}

We did not acquire an independent molecular-validation cohort (new IHC or sequencing);
PTEN and AR-activity's molecular concordance rests on already-public TCGA-PRAD multi-assay
data, and ERG's biological mechanism (why residual signal predicts grade at all) remains an
open, unverified hypothesis. LEOPARD's public labels carry no clinical covariates, so an
incremental-value-over-clinical-baseline claim is untestable with this cohort; the pool
comparison in Section~\ref{sec:leopard-results} is accordingly scoped to outcome association,
not incremental clinical value. Marker 7 is a post-hoc discovery, not a protocol-frozen
candidate, and its TCGA-PRAD outcome definition differs in granularity from LEOPARD's; among
censored TCGA-PRAD patients, marker 7's risk score correlates with follow-up duration more
than the corresponding LEOPARD check did ($\rho=-0.28$, $p=3\times10^{-5}$ vs.\ $p=0.07$ on
LEOPARD), a residual concern we could not fully resolve given TCGA-PRAD's multi-institution
heterogeneity and coarser outcome definition, though the landmark analysis argues against a
pure observation-time artifact. Two further limitations emerged from expanding this audit in
response to external review (Sections~\ref{sec:res-marker7-clinical}
and~\ref{sec:res-label-benchmark}). The completed R4 comparison uses the same 153-patient
complete-case cohort and shared patient-bootstrap draws for both endpoints. Multiple imputation
was not performed. The reconstructed grade-level delta C-index interval excludes zero, but its
paired IBS interval includes zero, the fuller-baseline and hierarchy intervals include zero,
and all official-PFI paired C-index and IBS intervals include zero. Thus ``grade-independent''
should not be read as ``independent of standard clinical covariates generally,'' and the result
does not establish robust independent or incremental prognostic value. Re-scoring marker 7
directly against TCGA's standardized PFS/DFS endpoints, rather than our
reconstructed label, gives a visibly attenuated (though still above-chance) transfer, indicating
the original headline C-index is partly specific to our label's construction. Official
TCGA-CDR PFI was evaluated directly (270 patients, 42 events; C-index 0.586, 95\% CI
0.482--0.689); PFS, DFS, and reconstructed recurrence are not treated as PFI substitutes.
What specific
morphological content the shared CONCH/Virchow recurrence-predictive direction reflects remains
unresolved: a preliminary, qualitative tile-level review ($n=6$ patients, 18 tiles, no
pathologist review) suggested loss of organized benign stromal architecture as a plausible
pattern, but one high-scoring tile showed an atypical color cast raising the possibility of a
staining or scanning artifact, and we present this only as hypothesis-generating. The
completed Gate A grid (Section~\ref{sec:failure-modes}) contains 360 correlated cells, not 360
independent tests or replications. Its five-seed Student-$t$ intervals quantify repeated tile
sampling on the same frozen patients/folds rather than patient or population uncertainty.
Marker-7 source retention varies from 498 to 502 patients; source-constant sensitivity requires
the completed common-498 R2 analysis. It preserved every audited contrast direction and had no
raw-to-common null crossing, but individual C-indices moved by up to 0.023. Because this is a
correlated post-hoc sensitivity on the same source and target cohorts, it neither identifies a
causal scale mechanism nor establishes universal encoder superiority or independent
replication. The official-PFI interval includes 0.5 and provides directionally
non-contradictory but not strong support. Gate A also records 60 coordinate shards and a historical
limitation: per-run command/config manifests were unavailable. A complete 463/463 NADT TIFF
header scan produced invalid-page-offset warnings for two slides. Those header findings do not
establish pixel-level impact. R5 closes the planned C04/C05 site, metadata, and power tables,
but not their scientific uncertainty: all six AR LOO intervals include zero; ScanScope is a raw
identifier available for 280/300 slides, not a scanner model; explicit stain metadata are
unavailable for all 300; and the SPOP MDE values are fixed-score planning approximations, not
effect-exclusion bounds. The completed R3/R4 analyses close the planned dependencies while
retaining endpoint-conditioned uncertainty.

\subsection{Conclusion}

A statistically significant association between a pathology foundation-model embedding and a
clinical or molecular label is a starting point, not an endpoint. Applied exhaustively and with
a prespecified protocol to six frozen candidates, and with the same audit tools applied
transparently to a seventh, post-hoc candidate, this kind of audit separates genuinely
transportable signal (grade, phenotype) from markers whose apparent strength did not survive a
stricter, nested re-analysis (PTEN, AR) and from configuration-sensitive evidence (SPOP).
The SPOP characterization establishes neither a robust positive nor a robust null: it is not a
biological-null or mechanism claim, and a modest effect and tumor dilution remain unresolved.
Its most consequential test shows that qualification for one target is not evidence of
qualification for another: a post-hoc marker fit directly against recurrence transferred
zero-shot to an independent cohort, while every prospectively qualified marker, combined,
failed to. The prespecified protocol, the discipline of re-testing a result under a stricter
design when a legitimate concern is raised, and the pitfalls this audit surfaced are offered as
reusable tools for the next study asked to make the same distinction.
